\part*{Part I --- Generation}
\addcontentsline{toc}{part}{Part I --- Generation}
\renewcommand{\thesection}{\arabic{section}}
\setcounter{section}{0}
\setcounter{definition}{0}

\bigskip
\noindent
\emph{This part establishes the generator. A presentation of a computational
system --- a signature of term constructors, a set of rewrite rules, and a choice
of how witnesses are collected --- determines a logic, and the determination is
mechanical for two of the three components. A fourth parameter, the algebra in
which observations take their values, is introduced and shown to be independent of
the other three. The part closes with a table exhibiting the logics used
throughout this document as settings of the four parameters. The material of
\S\ref{sec:realizability}--\S\ref{sec:dialintro} is a compression of
\cite{Zeus}; the fourth dial is from \cite{GradedWhere,WeightedGSLT}; the
adequacy discussion is from \cite{GsltOmnibus}.}
\bigskip

\section{Realizability, and why the seams matter}
\label{sec:realizability}

Most accounts take a logic as given: a fixed vocabulary of connectives, a notion
of model, and a truth predicate floating above computation, with truth values
arriving from a Boolean or Heyting algebra that is postulated rather than built.
For a discipline that wants its reasoning systems to be inspectable --- traceable
to constructions whose behaviour can be run and audited --- this is the wrong
starting point, because the glue binding truth to the algebra is opaque. Once it
is in place one can no longer see what the logic is \emph{made of}.

The commitment that makes the seams visible is realizability \cite{Kleene45,
vanOosten08}: a formula denotes the collection of computations that witness its
computational content. Under that commitment, and one further fixing --- of
\emph{how} witnesses are aggregated --- a logic factors into three components, and
every connective is of one of three kinds.

There is a further reason to care, beyond hygiene. If the connectives are
operations of three named kinds, then a new logic is obtained by turning dials
rather than by invention, and an algorithm defined over one logic can be
transported to another by transporting the dial settings. That is the difference
between a logic and a \emph{procedure for generating logics}, and it is the
practical content of this document. It also bears directly on a claim the AI
community makes on its own behalf: that important aspects of intelligence can be
faithfully rendered as computation. Logic has long been cast as the calculus of
reason. Exhibiting logic itself as \emph{arising from} computation is direct
support for that claim, since the calculus of reason is thereby shown to be, at
bottom, computational.

\begin{notation}
Throughout, $\T$ is the type of terms of the system under study, $\Coll$ is the
collection type in which witnesses are aggregated, and $\sem{\varphi} : \Coll(\T)$
is the denotation of $\varphi$. We write $e \real \varphi$ for ``$e$ witnesses
$\varphi$''.
\end{notation}

\section{The three components}
\label{sec:three}

\paragraph{Structural connectives.} Each term constructor lifts to a predicate
transformer. If $f : \T^{n} \to \T$ is an $n$-ary constructor, define
\[
  \sem{\widehat{f}(\varphi_{1},\dots,\varphi_{n})}
  \;=\;
  \{\, f(e_{1},\dots,e_{n}) : e_{i} \real \varphi_{i} \,\}.
\]
The lifted constructor $\widehat{f}$ speaks about the \emph{shape} of a witness:
it asserts that the witnessing computation was built by $f$ from sub-witnesses of
the $\varphi_{i}$. For the \rhoc{} calculus these are the connectives of namespace
logic \cite{MeredithRadestock05a,MeredithRadestock05b}: a connective for output,
one for input, the separating conjunction $\varphi \mid \psi$ coming from parallel
composition, and --- distinctively, and because the calculus is reflective --- the
name predicate $\quot{\varphi}$, which holds of a name whose quoted process
satisfies $\varphi$. Names see into structure.

\paragraph{Behavioural connectives.} Each rewrite is labelled by the context in
which it fires. A step $\Ktx[l] \to \Ktx[r]$ carries the one-hole context $\Ktx$,
and lifting labels to modalities gives
\[
  \sem{\langle \Ktx\rangle \varphi}
  = \{\, e : e \xrightarrow{\,\Ktx\,} e',\ e' \real \varphi \,\},
  \qquad
  \sem{[\Ktx]\varphi}
  = \{\, e : \forall e'.\ e \xrightarrow{\,\Ktx\,} e' \Rightarrow e' \real \varphi \,\},
\]
the context-indexed Hennessy--Milner modalities. There is one primitive modality
per rewrite rule \emph{and choice of redex position}; the quantifiers
$\exists j.\langle \Ktx_{j}\rangle \varphi$ and $\forall j.[\Ktx_{j}]\varphi$
range over redex positions.

The indexing is not bookkeeping. In the noughts-and-crosses development
\cite{GameNote} every move is the same action --- a square is claimed --- so an
action-labelled logic sees one label and cannot tell nine moves apart, while the
context-labelled logic sees as many labels as there are empty squares, and the
branching factor of $\exists j$ \emph{is} the legal move count. Getting this
right required a correction: the redex context must be instantiated, which forced
the encoding to make an empty square a receipt offering to be claimed, so that a
ply is exactly one rewrite.

\begin{remark}[The two halves of OSLF]
\label{rem:oslf}
Structural and behavioural connectives are the introduction and elimination
halves of an Operational-Semantics-in-Logical-Form presentation. They are
generated, respectively, from the constructors and from the labelled transitions,
and neither involves any choice once the presentation is fixed. This is the
precise content of the claim that the logics here are not \emph{ad hoc}.
\end{remark}

\paragraph{Collection connectives.} Any operation on $\Coll$ lifts to a
connective. With $\Coll = \Pfin$ the semilattice operations give the familiar
\[
  \sem{\varphi \wedge \psi} = \sem{\varphi} \cap \sem{\psi},
  \qquad
  \sem{\varphi \vee \psi} = \sem{\varphi} \cup \sem{\psi},
\]
but this is exactly where the choice of container is felt, and Part II is about
what varying it does.

\begin{principle}[Where the freedom is]
\label{prin:freedom}
Of the three components, two are determined by the presentation and one is a
parameter. Design freedom in a generated logic is concentrated in the collection
component --- and, once the fourth dial of \S\ref{sec:dialintro} is added, in the
value algebra.
\end{principle}

\section{What the rewrite rules see, and what adequacy costs}
\label{sec:adequacy}

Two consequences of Remark~\ref{rem:oslf} should be recorded before they cause
trouble later.

\paragraph{The rules determine the resolution of the modalities.} The generated
behavioural layer sees exactly what the rewrite relation distinguishes. If the
rules are presented as interleaving, the logic sees interleaving; if they are
presented so that independent steps commute as a genuine concurrency, the logic
sees true concurrency. There is no separate decision to be made about which
semantics the logic has. The presentation already made it.

\paragraph{Adequacy is a property of the target, not of the generator.} It is
tempting to say that a generated logic is automatically adequate --- that its
formulae separate precisely what interaction separates. This is false in general
and the error is worth naming, since an earlier statement of this programme made
it. Adequacy requires the target logic specification to supply enough structure,
in particular something like both a conjunction and a negation. A generator can
be asked to emit a logic without negation --- there are excellent reasons to,
discussed in Part IV --- and the resulting logic will not be adequate. The
Hypercube family of generators, which emits type systems rather than modal
logics, offers no adequacy theorem at all, and type systems are not generally
expected to be adequate \cite{GsltOmnibus}.

This is not a defect but a dial of its own, and it is the dial a practitioner
reaches for first. A user who wants model-checking to terminate gets it by
restricting the fragment of the system, or the set of admitted connectives, or
both. Restricting the system's fragment \emph{necessarily} restricts the
structural and behavioural connectives, since those are generated from it. When
model-checking terminates, the logic is necessarily not adequate for the whole of
a Turing-complete system. The methodological posture throughout is therefore:
\emph{the ideal logic defines; restricted fragments decide.} Where an infinitary
check is unavoidable, it is bounded by a token budget rather than forbidden ---
which, as Part IV shows, turns out to be the same budget that funds the learner's
observations.

\begin{remark}[A caveat about the structural layer]
\label{rem:spatialstrict}
The structural connectives are strictly more discriminating than
context-labelled bisimulation: the spatial layer can separate terms that no
interaction separates. This looks like a contradiction with adequacy and is not.
For interactive systems there is a uniform extension of the presentation ---
adjoin an atom for each constructor, overload the interaction constructor so that
applying it to that atom rebuilds the term, and add the corresponding destructors
--- which makes every structural distinction visible to bisimulation, at which
point the two notions agree. The construction is folklore in the spatial-logic
community, following Caires' observation of the phenomenon \cite{Caires04,
CairesCardelli03}, and is made explicit in a companion note. We use the fact and
do not reprove it.
\end{remark}

\section{The fourth dial: what an observation returns}
\label{sec:dialintro}

The three components above answer \emph{which} connectives a logic has. They do
not answer what a formula, applied to a term, returns. The default answer is a
truth value in $\B = \{\bot,\top\}$, and it is so thoroughly the default that it
is rarely recorded as a choice.

Record it. Let $\V$ be an algebra of values, and let $\sem{\varphi}(P) \in \V$.
For the positive fragment, a commutative semiring suffices: conjunction is
multiplication, disjunction is addition, and the constants $0,1$ are the units.
The instances used in this document are collected in Table~\ref{tab:valuealgebras}.

\begin{table}[h]
\centering
\renewcommand{\arraystretch}{1.25}
\begin{tabular}{@{}llll@{}}
\toprule
$\V$ & $\oplus$ & $\otimes$ & reading \\
\midrule
$\B$                & $\vee$ & $\wedge$ & crisp: the classical guard \\
$[0,1]$             & $\max$ & $\times$ & possibility; probabilistic logic \\
$\R_{\ge 0}$        & $+$    & $\times$ & propensity; stochastic simulation \\
Viterbi $([0,1])$   & $\max$ & $\times$ & best explanation, with a witness \\
tropical $(\R\cup\{\infty\})$ & $\min$ & $+$ & least cost, with a witness \\
$\Cx$               & $+$    & $\times$ & amplitude; interference \\
a quantale          & $\bigvee$ & $\otimes$ & graded truth with fixed points \\
\bottomrule
\end{tabular}
\caption{Value algebras used in this document. The last row is developed in
Part II, where it turns out to be the same object as the container of classical
linear logic.}
\label{tab:valuealgebras}
\end{table}

Three observations about this dial, each developed later.

\paragraph{It is independent of the other three.} The connectives available are
fixed by $(\Sig,\Rw,\Coll)$; what they return is fixed by $\V$. One can run the
same spatial-behavioural formula over $\B$, over $\R_{\ge 0}$ and over $\Cx$ and
get, respectively, a model-checking question, a propensity, and an amplitude.

\paragraph{It is where negation goes to die.} The positive fragment needs only a
semiring, but negation needs more. On $[0,1]$ the map $1-x$ serves; on $\Cx$
there is no canonical candidate and De Morgan fails. The design adopted in
\cite{GradedWhere} is therefore two-sorted: a \emph{crisp} sort with full Boolean
structure and fixed points, and a \emph{graded} sort with semiring structure and
neither negation nor fixed points. Part II explains \emph{why} that is the right
line to draw, by identifying the value algebras in which negation is available.

\paragraph{It changes what the logic is for.} Over $\B$ a formula is a question.
Over $\R_{\ge 0}$ it is a rate. Over $\Cx$ it is a contribution to a sum over
histories, and disjunction acquires the possibility of being false with both
disjuncts true --- which is to say that interference is not an exotic addition to
the logic but the ordinary behaviour of $\oplus$ once additive inverses exist.

\section{In what sense the logics are not \emph{ad hoc}}
\label{sec:notadhoc}

It is worth being exact about the claim, since it is easy to overstate.

\paragraph{What is claimed.} Given a presentation $(\Sig, \Rw)$ and a choice of
$(\Coll, \V)$, the connectives of the logic and their semantics are determined.
There is no step at which a designer says ``and let us also have a connective for
$\dots$''. Two people who fix the same four parameters write down the same logic.

\paragraph{What is not claimed.} That the four parameters are themselves forced.
They are not; choosing them is the whole art. Nor that the resulting logic is
\emph{good} --- adequate, decidable, or well-behaved. Nor that every logic in the
literature arises this way. The claim is about the map from parameters to logics,
not about its surjectivity.

\paragraph{Why this is more than tidiness.} Three dividends recur.

First, \emph{transport}. An algorithm written against the generated logic of one
system transports to another by transporting the parameters. The knot classifier,
the finite-geometry encoding, the board-game ladder and the swarm model are the
same construction at four settings.

Second, \emph{prediction}. Because the connectives are forced, one can predict
which connectives a given setting will and will not have, and be surprised when a
system in the wild turns out to have exactly those. The semitopology case of
Part II is the sharpest instance: the framework predicts, from closure properties
of a container alone, that the algebraic dual of a semitopology cannot have a
meet and must instead carry a symmetric compatibility relation --- which is
exactly the structure Gabbay arrived at independently \cite{GabbaySemiframes}.

Third, \emph{accountability}. When the logic is wrong for a purpose, one can say
which dial was set wrong, rather than starting over.

\section{The instantiation table}
\label{sec:table}

Table~\ref{tab:instances} exhibits the logics used in this document as settings of
the four parameters. It is the central claim of Part I in one object, and the rest
of the document is an unpacking of its rows.

\begin{table}[h]
\centering
\small
\renewcommand{\arraystretch}{1.3}
\begin{tabular}{@{}>{\raggedright\arraybackslash}p{3.1cm}>{\raggedright\arraybackslash}p{2.6cm}>{\raggedright\arraybackslash}p{2.9cm}>{\raggedright\arraybackslash}p{1.6cm}>{\raggedright\arraybackslash}p{3.4cm}@{}}
\toprule
logic & $\Sig,\Rw$ & $\Coll$ & $\V$ & where it appears \\
\midrule
namespace / spatial-behavioural logic
  & \rhoc{} calculus
  & $\Pfin$
  & $\B$
  & senses and assays \cite{ScientistNote}; the classifier notes \\
intuitionistic (geometric)
  & any
  & frame of opens
  & $\B$
  & the ambient reading of affirmation \cite{Vickers89} \\
classical linear logic
  & any
  & $\Pfin(\Mset\,A)$, with a pole
  & $\B$
  & \S\ref{sec:linear}; resource discipline \cite{Girard87} \\
Gabbay's coalition logic
  & participants
  & $\Open(P)$, a semiframe
  & $\B$
  & consensus \cite{Gabbay25,ConsensusTypes} \\
the graded coalition modality
  & \rhoc{} calculus
  & $\Open(P)$
  & tropical / Viterbi
  & herd and swarm \cite{QuorumNote} \\
graded where-clauses
  & \rhoc{} calculus
  & $\Pfin$
  & $[0,1]$, $\R_{\ge0}$, $\Cx$
  & the instrument \cite{GradedWhere} \\
quantale-graded \rhoc{}
  & \rhoc{} calculus
  & $\Pfin$
  & a quantale
  & fixed points survive \cite{FuzzyRho} \\
the choice-hierarchy formulae
  & \rhoc{} calculus
  & $\Mset$ of candidates
  & any of the above
  & altitude and price \cite{ChoiceFormulae} \\
weighted rewriting
  & any MeTTaIL GSLT
  & $\Pfin$
  & $\R_{\ge0}$, $\Cx$
  & simulation \cite{WeightedGSLT} \\
\bottomrule
\end{tabular}
\caption{Logics as settings of the four parameters.}
\label{tab:instances}
\end{table}

Two rows of that table deserve immediate comment, because they are the ones a
reader is most likely to think cannot be right.

The linear logic row is worked in \S\ref{sec:linear}. The claim is not that linear
logic can be \emph{modelled} in this framework --- of course it can, it is a logic
--- but that it is what the generator emits at a particular container setting, and
that its two conjunctions, its exponential and its involutive negation each
correspond to an identifiable feature of that container.

The coalition row is the one that comes with independent corroboration. It is
also the row that discharges a promise: \cite{Zeus} treats only the collection
component of coalition logic and defers the structural and behavioural components,
together with the identification of the underlying points with \rhoc{} terms, to a
companion. Part II, \S\ref{sec:semitop}, is that companion.
